\section{Sharp translation regimes}
\label{sec:regime-criteria}

\begin{theorem}[Exact terminal criterion and regime corollaries]
\label{thm:regime-criterion}
Consider a family indexed by \(X\), with anchors \(r_X\), lengths
\(N_X\to\infty\), primitives \(L_X\), and ordinary sums \(S_X\).
Assume
\[
 \cT_{0,N_X}(L_X)=o(1).
\]
Then
\begin{equation}
 \frac{S_X}{N_X}=o(1)
 \quad\Longleftrightarrow\quad
 \frac{r_X}{N_X}L_X(N_X)=o(1).
\label{eq:exact-terminal-criterion}
\end{equation}
Consequently:
\begin{enumerate}[label=\textup{(\alph*)},leftmargin=2.5em]
\item If \(r_X=0\), no terminal condition is needed.
\item If \(0<r_X=o(N_X)\), the exact terminal scale is
  \(L_X(N_X)=o(N_X/r_X)\); the stronger bound
  \(L_X(N_X)=O(1)\) is a convenient sufficient condition.
\item If \(r_X\asymp N_X\), then the two conditions
  \[
  \cT_{0,N_X}(L_X)=o(1),\qquad L_X(N_X)=o(1)
  \]
  imply \(S_X/N_X=o(1)\).
\item If \(r_X/N_X\to\infty\), then
  \[
  \cT_{0,N_X}(L_X)=o(1),\qquad
  L_X(N_X)=o(N_X/r_X)
  \]
  imply \(S_X/N_X=o(1)\).
\end{enumerate}
The terminal condition in \eqref{eq:exact-terminal-criterion} is
necessary and sufficient without additional structure.
\end{theorem}

\begin{proof}
The identity
\[
 \cT_{r_X,N_X}(L_X)
 =
 \cT_{0,N_X}(L_X)
 +\frac{r_X}{N_X}L_X(N_X)
\]
and
\(S_X/N_X=\cT_{r_X,N_X}(L_X)\) prove
\eqref{eq:exact-terminal-criterion}.  The regime statements are
immediate.  The examples in
\cref{ex:constant-obstruction,ex:tuned-anchor} show that the anchor
term cannot be removed.
\end{proof}

For a positive comparison scale define the translated harmonic mass
\begin{equation}
 \cH_{r,N}:=\sum_{n=1}^N\frac1{r+n}.
\label{eq:harmonic-mass}
\end{equation}

\begin{lemma}[Short translated harmonic mass]
\label{lem:short-harmonic-mass}
For \(r\ge N\ge1\),
\begin{equation}
 \frac{N}{r+N}\le\cH_{r,N}\le\frac Nr,
 \qquad
 \left(1+\frac rN\right)\cH_{r,N}\le2.
\label{eq:short-harmonic-bounds}
\end{equation}
\end{lemma}

\begin{proof}
Each denominator lies between \(r+1\) and \(r+N\).  The displayed
bounds follow by comparison; the last uses \(N/r\le1\).
\end{proof}

\begin{theorem}[Prefix-uniform criterion for short translated fibers]
\label{thm:short-prefix-uniform}
Assume \(r\ge N\) and
\begin{equation}
 \max_{1\le m\le N}
 \frac{|L_{r,m}|}{\cH_{r,m}}
 \le\varepsilon.
\label{eq:relative-prefix-input}
\end{equation}
Then
\begin{equation}
 \left|\frac{S_N}{N}\right|\le3\varepsilon.
\label{eq:short-prefix-output}
\end{equation}
\end{theorem}

\begin{proof}
By \cref{thm:anchored-identity},
\begin{align*}
 \left|\frac{S_N}{N}\right|
 &\le
 \left(1+\frac rN\right)|L_{r,N}|
 +\frac1N\sum_{m<N}|L_{r,m}|\\
 &\le
 2\varepsilon+
 \frac{\varepsilon}{N}\sum_{m<N}\cH_{r,m}.
\end{align*}
Since \(\cH_{r,m}\le m/r\le1\), the last average is at most
\(\varepsilon\).
\end{proof}

\begin{remark}[The quantifier is strong]
\label{rem:prefix-quantifier}
The hypothesis is uniform in every prefix and normalized by its
actual translated harmonic mass.  A statement only at the terminal
endpoint, or one with an \(o(1)\) error not measured relative to
\(\cH_{r,m}\), does not imply \eqref{eq:relative-prefix-input}.
Moreover the case \(m=1\) reads \(|a_1|\le\varepsilon\).  For a
unit-modulus arithmetic sequence this cannot hold with
\(\varepsilon=o(1)\).  A realistic application must peel a finite
initial segment, return it exactly to the outer sum, and impose the
relative estimate only on the remaining prefixes.  The theorem is a
strong sufficient certificate, not a claim that literal M\"obius
fibers already satisfy it.
\end{remark}

\begin{corollary}[Cutoff-prefix certificate]
\label{cor:cutoff-prefix}
Assume \(r\ge N\), fix \(1\le M\le N\), and suppose
\[
 \sup_{M\le m\le N}
 \frac{|L_{r,m}|}{\cH_{r,m}}
 \le\varepsilon.
\]
Then
\begin{equation}
 \left|\frac{S_N}{N}\right|
 \le
 3\varepsilon
 +
 \frac1N\sum_{m=0}^{M-1}|L_{r,m}|.
\label{eq:cutoff-prefix-output}
\end{equation}
\end{corollary}

\begin{proof}
Use \cref{eq:anchored-defect-identity}.  The terminal term is at
most
\[
 \left(1+\frac rN\right)\varepsilon\cH_{r,N}
 \le
 \varepsilon\left(1+\frac Nr\right)
 \le2\varepsilon.
\]
The controlled part of the backward average is at most
\[
 \frac{\varepsilon}{N}
 \sum_{m=M}^{N-1}\cH_{r,m}
 \le\varepsilon,
\]
and the omitted prefixes give the last term in
\eqref{eq:cutoff-prefix-output}.
\end{proof}

\begin{proposition}[Smooth terminal bound]
\label{prop:smooth-terminal}
Suppose, for every \(0\le m\le N\),
\[
 L_{r,m}=\ell_{r,m}+e_{r,m},
\]
where
\[
 \cT_{0,N}(\ell)=0.
\]
Then
\[
 \left|\frac{S_N}{N}\right|
 \le
 \frac rN|\ell_{r,N}|
 +\left(1+\frac rN\right)|e_{r,N}|
 +\frac1N\sum_{m<N}|e_{r,m}|.
\]
Thus a model primitive transfers only after its anchor term and the
full prefix error have both been controlled.
\end{proposition}

\begin{proof}
Insert \(L=\ell+e\) into
\eqref{eq:anchored-defect-identity} and use the assumed cancellation
of the unshifted model defect.
\end{proof}
